\chapter{Experimental Evaluation}
\label{ch:evaluation}
\chaptermark{Experimental Evaluation}

We test construction/capacity (RQ1), execution (RQ2), budget quality/reuse (RQ3) and competitive regions (RQ4).

\section{Experimental Protocol}
\label{sec:7-1}

\subsection{Comparison protocol and reproducibility}

\Cref{sec:3-2,sec:3-3,sec:3-5} define the protocol. Search excludes construction/calibration; each ablation has its own baseline.

CPU JHQ uses unmatched settings/protocols: October Vogue has 1,024 lists and $[8,4]$ bits versus E1's 4,096 lists. A matched CPU/GPU speedup is unavailable.

Adaptive calls/device-only timing are in \cref{sec:a-e1-protocol}.

Metrics and frontier estimators follow \cref{sec:3-2,sec:3-4};
$\glut$ and $\gscan$ retain the distinct meanings in \cref{sec:5-1}.





\section{Modern GPU ANN Recall--QPS Frontiers}
\label{sec:7-2}

\subsection{Comparison scope and frontier construction}

\Cref{fig:7-frontiers} compares JHQ-GPU with the graph, compressed-IVF
and brute-force baselines of \cref{sec:3-4} \cite{cagra,pq,ivfrabitq},
using the first four datasets in \cref{tab:3-data}. Each uses one batch
of 1,000 queries with common GT. T5 separately evaluates small batches
in \cref{sec:7-smallbatch}; \cref{sec:7-6} reports internal ablations.


The recorded schedule uses $\gscan=0$ for all 105 arXiv/OpenAI
run--budget cells and $\gscan=4$ for all 27 Vogue cells, triggered by
Vogue's candidate-count CV of 0.338--0.622. Thus the non-Vogue frontiers
exercise query-centric scanning, while Vogue also exercises grouped
retention, shared bounds and merging (\cref{sec:5-5}).

Applying the frontier rule in \cref{sec:3-4}, JHQ contributes 27--39
points per dataset and each baseline family 54--264. E1 contains 1,590
valid points, no failure records and no invalid-ID or duplicate-row flags.
These points share indexes across parameter sweeps.

Grids are in \cref{tab:3-grids,sec:b-grids}. A curve ending short of a target has no tested coverage there.


\begin{figure}[tbp]
\centering
\includegraphics[width=\linewidth]{ch7/frontiers.pdf}
\caption[E1 recall--throughput frontiers]{Measured host-to-host recall--QPS frontiers on the four E1 datasets.
The displayed recall window begins at 0.90, with a vertical line at 0.97;
all baseline families,
including CAGRA int8, participate in frontier construction. Lines connect
measured points; throughput uses a logarithmic axis.}
\label{fig:7-frontiers}
\end{figure}

Nine-bit RaBitQ shares the nominal $9d/8$ payload; 4/8-bit codes use less, so storage is unmatched.
\Cref{tab:4-bytes,sec:7-7} give representation costs and measured
allocations alongside this performance envelope.


\subsection{Matched recall and opposing operating regions}

\Cref{tab:7-matched} reports recall 0.97 using the interpolated and
qualifying-measured estimators defined in \cref{sec:3-4}.


\begin{table}[tbp]
\centering
\caption[Cross-system comparison at recall 0.97]{Recall 0.97 comparison. Each ratio is interpolated / measured;
values above one favour JHQ. JHQ kQPS is the interpolated estimate.
The small arXiv--RaBitQ difference changes direction with the estimator.}
\label{tab:7-matched}
\small\setlength{\tabcolsep}{4pt}
% Generated by scripts/build_ch7.py; do not edit numbers manually.
\begin{tabular}{@{}lrrr@{}}
\toprule
Dataset & JHQ kQPS & vs CAGRA fp32 & vs IVF-RaBitQ \\
\midrule
Vogue & 230.8 & 1.416 / 1.537 & 1.059 / 1.046 \\
arXiv & 258.6 & 1.317 / 1.422 & 0.996 / 1.015 \\
OpenAI-1536 & 126.7 & 1.747 / 2.294 & 1.130 / 1.111 \\
OpenAI-3072 & 55.9 & 1.747 / 2.088 & 1.157 / 1.346 \\
\bottomrule
\end{tabular}

\end{table}

At batch 1,000/recall 0.97, removing grouping drops Vogue's JHQ/RaBitQ ratio from 1.059 to 0.933 at fixed index, probes, budget and final outputs. This is an integrated grouping benefit; the bound alone is not isolated.

CAGRA int8 leads by $1.20$--$1.98\times$ at recall 0.90--0.92 but lacks 0.97 coverage. RaBitQ wins some 0.98+ targets; CAGRA fp32 wins near 0.99 on arXiv (\cref{fig:7-frontiers}).

\subsection{Precision sensitivity on the OpenAI datasets}
\label{sec:7-precision}

Supplementary source inspection establishes that the E1 implementation
uses TF32 multiplication with FP32 accumulation for OpenAI rotation and
centroid GEMMs, despite FP32 input/storage and a misleading source comment
(\cref{sec:4-2}). The main E1 frontier remains the historical measurement
of that implementation. In campaign D, a separate FP32 variant changes the cuBLAS
math-mode switch and replays 37/39 OpenAI configurations at batch 1,000,
five rounds each. Execution order alternates between variants across
probe settings. Both reuse the same TF32-trained model; the intervention
changes add/query arithmetic, including database encoding and routing.
It is not a comparison of independently FP32-trained models.
\Cref{tab:7-precision} summarises the paired changes.

\par\medskip\noindent\begin{minipage}{\linewidth}
\centering\captionsetup{hypcap=false}
\captionof{table}[FP32 sensitivity at matched configurations]{FP32 sensitivity. QPS ratio is FP32/TF32 at equal configuration
(geometric mean); recall differences have the same direction. The last
two columns compare the FP32 frontier at recall 0.97 with archived E1
IVF-RaBitQ and CAGRA fp32 frontiers. Campaign D does not remeasure the baselines.}
\label{tab:7-precision}
\scriptsize\setlength{\tabcolsep}{3pt}
\begin{tabular}{@{}lrrrrr@{}}
\toprule
Dataset & Pairs & QPS ratio & $\Delta$Recall & vs RaBitQ & vs CAGRA \\
\midrule
OpenAI-1536 & 37 & 0.953 & +0.0003--+0.0005 & 1.086 & 1.679 \\
OpenAI-3072 & 39 & 0.951 & +0.0000--+0.0003 & 1.079 & 1.629 \\
\bottomrule
\end{tabular}

\end{minipage}\par\medskip

At equal configuration, FP32 reduces throughput by about 5\% and increases
recall by 0--0.0005. At recall 0.97, FP32 is about 3\% slower than same-session TF32, yet $1.086/1.079\times$ archived RaBitQ and $1.679/1.629\times$ archived CAGRA fp32. New-session v69 gives 1.118/1.114 against archived RaBitQ, versus historical 1.130/1.157. D compares baselines across sessions; only v69/v72 share a session. \Cref{fig:7-precision}
shows the two JHQ frontiers with their archived comparators.
T1 separately remeasures same-session anchors, with grid-sensitive results in \cref{sec:b-precision-anchors}.

\begin{figure}[tbp]
\centering
\includegraphics[width=\linewidth]{supplement/precision.pdf}
\caption[TF32 and FP32 precision frontiers]{Campaign D precision frontiers. Solid curves are same-session
JHQ variants; dashed/dotted baselines are historical E1 measurements.
All use the chapter's frontier and interpolation rules.}
\label{fig:7-precision}
\end{figure}


\subsection{Interpretation and comparison limits}

Both OpenAI datasets use Qdrant DBpedia embeddings \cite{data_openai1536,data_openai3072}. After truncation/renormalisation, 1,000 query and 10,299 sampled base rows have cosine $\geq0.9980$ (median 1.000000), with GT overlap 9.046/10. This supports row alignment without entity IDs. Doubling width leaves the CAGRA ratio at 1.747 and changes RaBitQ's 1.130 to 1.157: width alone does not explain the gain.

Gains accompany smaller 128-probe anchors: OpenAI 4--8, Vogue 64, arXiv 100 (later FULL 200). This is an association, not an isolated effect or matched-frontier setting.

\Cref{sec:7-8} checks sensitivity to small margins and the interface/memory/compilation differences of \cref{sec:3-3}.


\subsection{Small batches change the competitive region}
\label{sec:7-smallbatch}

T5 compares scheduled JHQ, CAGRA fp32/int8 and IVF-RaBitQ at batches
8/32 on the four E1 datasets. Five synchronous host-to-host rounds use
1,000/992 queries, respectively. Its 1,808 points include JHQ's one-block
control; IVF-PQ and brute force are outside this campaign.
\Cref{fig:7-smallbatch,tab:7-smallbatch} compare the resulting regions.

\suppressfloats[t]
\begin{figure}[htbp]
\centering
\includegraphics[width=\linewidth]{t4-t5/smallbatch.pdf}
\caption[Small-batch competitive regions]{T5 JHQ / fastest-competitor
QPS under interpolation with qualifying measured endpoints. Above one
favours JHQ. \Cref{tab:7-smallbatch} gives strict thresholds.}
\label{fig:7-smallbatch}
\end{figure}

\begin{table}[htbp]
\centering
\caption[Small-batch matched-recall ratios]{T5 QPS ratios:
interpolation with qualifying endpoints / strict measured threshold.
The fastest competitor is selected independently for each estimator;
higher-recall qualifying endpoints remain eligible without extrapolation.}
\label{tab:7-smallbatch}
\footnotesize\setlength{\tabcolsep}{4pt}
\begin{tabular}{@{}lrrrr@{}}
\toprule
Dataset & $b$ & $R=0.90$ & $R=0.95$ & $R=0.97$ \\
\midrule
Vogue & 8 & 1.155 / 1.091 & 1.073 / 0.980 & 0.893 / 0.854 \\
Vogue & 32 & 1.069 / 0.997 & 0.956 / 0.895 & 0.865 / 0.839 \\
arXiv & 8 & 1.479 / 1.454 & 0.859 / 0.851 & 0.697 / 0.698 \\
arXiv & 32 & 1.166 / 1.148 & 0.667 / 0.661 & 0.622 / 0.582 \\
OpenAI-1536 & 8 & 1.316 / 1.131 & 0.851 / 0.825 & 0.641 / 0.620 \\
OpenAI-1536 & 32 & 1.071 / 0.915 & 0.770 / 0.752 & 0.937 / 0.914 \\
OpenAI-3072 & 8 & 1.032 / 0.953 & 0.798 / 0.798 & 0.721 / 0.702 \\
OpenAI-3072 & 32 & 0.866 / 0.734 & 0.608 / 0.771 & 1.008 / 1.351 \\
\bottomrule
\end{tabular}

\end{table}

At recall 0.97 and batch 8, CAGRA fp32 leads on all four datasets under
both estimators: interpolated JHQ ratios are 0.641--0.893. Batch-32 ratios are 0.622--1.008: OpenAI-3072 is near parity under interpolation, JHQ-favourable under the measured threshold. Vogue's batch-8/0.95 lead also depends on interpolation. Borderline rankings depend on grid resolution.

Qualifying endpoints above the target remain eligible. At recall 0.90, batch-8/32 ratios are 1.032--1.479/0.866--1.166, with three/one strict-threshold JHQ wins. Dropping endpoints overstates leads. Scheduling gains $1.7$--$3.3\times$ over one block at 0.95, without ensuring cross-system wins.

Null-call medians are 0.012--0.072\,ms and remain included. CAGRA uses Python, JHQ/RaBitQ C++; warm-ups differ (\cref{sec:a-t4-t5}).

\section{Cartesian LUT, Packed Loads and Lane Refinement}
\label{sec:7-3}

The RQ2 ablations measure both kernel and application effects.

\subsection{Exact table factorization and its throughput consequence}

The LUT ablation tests the cost model in \cref{sec:5-2} and FP32 ranking
effects.

\Cref{fig:7-lut} compares Cartesian table grouping on six datasets,
probe counts 8/32/128/512,
fixed $\alpha=100$, and $\glut\in\{1,2,4,8\}$. Five rounds use batch 1,000 (Stella 999). Two groups lead in 22/24 cells. Exceptions are near ties: arXiv/128 favours one group by 0.23\%; OpenAI-3072/8 favours four by 0.84\%. Paired two-versus-one wins are 23/24; all dataset geometric means exceed one (1.021--1.318). Shared-index cells are not independent trials.

\begin{figure}[tbp]
\centering
\includegraphics[width=\linewidth]{ch7/lutgroups.pdf}
\caption[Cartesian lookup-table grouping]{Cartesian table grouping at fixed $\alpha=100$ (v71).
Group counts describe $\glut$, not parallel scan groups. Curves are medians
of per-dataset QPS ratios to the one-group arm; shading shows their range.}
\label{fig:7-lut}
\end{figure}

The two-group/one-group ratio ranges from $0.998$ to $1.914$.
Four and eight groups each store 16 entries but require more lookups
and additions. These arms show that smaller tables need not improve
throughput; occupancy and instruction costs were not isolated. The maximum
reported recall spread within any four-arm cell is $10^{-4}$.
At 128 probes and $\alpha=100$, T4 compares identical candidates on
Vogue, arXiv and OpenAI-3072. Full versus $16+16$ tables differ in
22.6--26.1\% of candidate-score bit patterns, with maximum relative error
$4.53\times10^{-7}$. Across all forms, 20.5--29.4\% differ, with maximum
relative error $4.64\times10^{-7}$. Vogue has 2--4/256 final-ID
differences, while arXiv/OpenAI have none. Unchanged primary survivor
sets can still produce different final IDs because the rounded primary
score enters refinement (\cref{sec:a-t4-t5}).


\subsection{Packed access and cooperative residual scoring}

Packing wins $1.012$--$1.658\times$ across 48 matched pairs (six datasets, four probe counts, $\glut\in\{1,2\}$), with recall differences at most $10^{-4}$. Both layouts are subspace-major. This tests packing/load/unpack together; one aggregate rate per arm gives less timing evidence than v71's five rounds.

\Cref{fig:7-packing-refinement,tab:7-refinement} then separate wider loads
from lane cooperation in residual refinement. L0 is the original scalar
path; L1 changes the load path without distributing a candidate across
multiple lanes; L16 assigns 16 lanes to a candidate. The same process
rotates arm order across five timing rounds, with three search calls per
arm per round. For the twelve ordinary configurations, L16 improves
end-to-end throughput by $1.015$--$1.789\times$. An additional arXiv
configuration at $\alpha=200$, explicitly a high-budget control, reaches
$2.383\times$; it is not the gain range of the ordinary selected budgets.

\begin{figure}[tbp]
\centering
\includegraphics[width=\linewidth]{ch7/packing_refinement.pdf}
\caption[Packed access and refinement controls]{Packing and refinement controls. Left: median packed-load QPS
 gain across six datasets, with shaded minimum--maximum ranges, for each
 table grouping and probe count. Right: measured load-only and 16-lane
 end-to-end gains at 128 probes and the ordinary dataset-specific budgets.
 Panels use different experiment families.}
\label{fig:7-packing-refinement}
\end{figure}

\begin{table}[tbp]
\centering
\caption[Refinement gains at 128 probes]{Refinement at 128 probes. Ratios are paired end-to-end gains over
 L0, except the last column, which is the profiled refinement-kernel gain.
 The kernel and application gains must not be interchanged.}
\label{tab:7-refinement}
\small\setlength{\tabcolsep}{4pt}
% Generated by scripts/build_ch7.py; do not edit numbers manually.
\begin{tabular}{@{}lrrrr@{}}
\toprule
Dataset & $\alpha$ & $L1/L0$ QPS & $L16/L0$ QPS & Kernel gain \\
\midrule
arXiv & 100 & 1.368 & 1.789 & 5.97 \\
BGE & 32 & 1.026 & 1.062 & 1.97 \\
OpenAI-1536 & 8 & 1.055 & 1.145 & 3.34 \\
OpenAI-3072 & 4 & 1.063 & 1.171 & 8.81 \\
Stella & 16 & 1.030 & 1.053 & 2.66 \\
Vogue & 64 & 1.161 & 1.240 & 3.72 \\
\bottomrule
\end{tabular}

\end{table}

The load-only control is bit-identical to L0 in the tested cases, yet
already improves arXiv throughput by $1.368\times$. Lane cooperation adds
a larger reduction in refinement time. On arXiv at $\alpha=100$, the
profiled refinement kernel falls from 4,277.5 to 716.8\,$\mu$s, a
$5.97\times$ gain, while total throughput improves $1.789\times$.
OpenAI-3072 gains $8.81\times$ in refinement and $1.171\times$ overall; routing, screening, selection and transfers remain.

For baseline wall-time fraction $f$ and refinement speedup $s$, the replacement model gives
\[
 \frac{T'}T=(1-f)+\frac f s,\qquad
 \frac T{T'}=\frac1{1-f+f/s}.
\]
ArXiv/OpenAI-3072 $f=0.529/0.168$ predicts $1.787/1.175\times$, versus measured $1.789/1.171\times$. Separate profiling makes this explanatory, rather than a paired timing identity.

Lane FP32 reassociation reaches relative distance error $\approx1.6\times10^{-6}$ and occasionally changes order. Sanitizer coverage is limited to recorded subsets; general bit identity is not claimed.

\section{Global Top-c, Shared Bound and Probe Correctness/Ablation}
\label{sec:7-4}




\subsection{Strict CPU reference for computed-score selection}
\label{sec:7-strict-selection}

T4 sorts GPU-produced scores by $(\primarydist,\mathrm{position})$ on CPU. The exporter follows scan layout/accumulation and matches valid returned scores bitwise; source comparison finds only export changes from the table-sweep version. Its 270 configurations cover Vogue/arXiv/OpenAI-3072, probes 1/8/128, $\alpha\in\{4,100,200\}$, $\gscan\in\{0,1,2,4,8\}$ and both bound settings: 69,120 instances (\cref{tab:7-strict-selection}).

\suppressfloats[t]
\begin{table}[htbp]
\centering
\caption[Strict computed-score selection checks]{T4 CPU-reference checks.
Set and order equality use the complete primary key. Valid count checks
$\min(c,n_c)$ survivors; final equality sorts retained computed refined
scores by $(\refineddist,\mathrm{ID})$, including their FP32 bits. Counts
are query instances, with queries reused across configurations.}
\label{tab:7-strict-selection}
\footnotesize\setlength{\tabcolsep}{3pt}
\begin{tabular}{@{}lrrrrr@{}}
\toprule
Path & Queries & Set / order exact & Tie-only & Valid count & Final exact \\
\midrule
$\gscan\geq1$ & 55,296 & 55,296 / 55,296 & 0 & 55,296 & 55,296 \\
$\gscan=0$ & 13,824 & 13,768 / 8,758 & 56 & 13,824 & 13,824 \\
\bottomrule
\end{tabular}

\end{table}

Grouped order matches in 55,296/55,296 instances. One-block distance order passes in 13,824/13,824, with 56 cutoff-tie set differences. All 19,120 underfull counts and 69,120 retained-score final selections pass. These test \cref{lem:5-buffer,prop:5-local-coverage,prop:5-shared-bound}; separate scan/bound and sanitizer checks have their own coverage (\cref{sec:a-t4-t5}).

\subsection{Shared bounds preserve the selection target}

\Cref{prop:5-shared-bound} establishes the selection guarantee. The
paired experiment measures the work saved by sharing that bound.

All 48 primary/final comparisons agree; throughput ratios are $0.994$--$1.562\times$. At arXiv $\alpha=200$, batch 1,000 and eight groups, sharing cuts merge 2,629$\to$471\,$\mu$s and scan 7,968$\to$5,513\,$\mu$s. Some cells regress slightly.

The pilot's first group holds 46--81\% of winners. Equal quotas need fallback; unequal ones leave a dominant group. Shared bounds avoid fixed quotas.

\subsection{Exact probe selection and the limits of pipeline identity}

The radix selector retains score arithmetic, including OpenAI TF32 (\cref{sec:7-precision}). All 1,792 synthetic cases pass; probes and primary positions/distances match on all 1,000 queries of arXiv, Vogue and OpenAI-3072.

\begin{table}[tbp]
\centering
\caption[Probe-selection throughput ratios]{v69/v65 end-to-end throughput ratios for the exact computed-score probe
 change. Each workload--batch comparison has ten timing rounds per binary,
 collected in alternating runs. Ratios $>1$ favour v69.}
\label{tab:7-probe}
\small
% Generated by scripts/build_ch7.py; do not edit numbers manually.
\begin{tabular}{@{}lrrrr@{}}
\toprule
Dataset & Batch 8 & 32 & 256 & 1,000 \\
\midrule
arXiv & 1.074 & 1.063 & 1.035 & 1.034 \\
Vogue & 1.109 & 1.103 & 1.052 & 1.041 \\
OpenAI-3072 & 1.167 & 1.117 & 1.055 & 1.050 \\
\bottomrule
\end{tabular}

\end{table}

\Cref{tab:7-probe} shows end-to-end improvements of $1.034$--$1.167\times$.
At batch 8, the profiled probe stage falls from 56.5 to 22.6\,$\mu$s
on arXiv, 52.7 to 14.7\,$\mu$s on Vogue, and 77.8 to 16.7\,$\mu$s
on OpenAI-3072. Unchanged stages limit full-batch gains. These ratios use their own baselines.

Vogue's one-group/batch-8 probe/primary dumps match bytewise across binaries on all 1,000 queries, yet final recall is 0.9642/0.9644. The baseline itself gives 0.9642 in the identity run and 0.9644 in both timing runs. In that controlled case, variation arises after primary selection and without changing probes; the cause is unresolved. Scheduled/one-block outputs match within each checked process (\cref{app:reproducibility}).

Optional v70 FP16 matches v69 probe sets on 936/1,000 arXiv, 979/1,000 Vogue and 998/1,000 OpenAI-3072 queries; mean overlap is high. High overlap and small recall changes do not make the
primitive exact. Those results are excluded from the main E1 frontier
and from the exact-probe claim.

\section{Held-out Budget Quality, Throughput and Payback}
\label{sec:7-5}

RQ3 compares the policy of \cref{ch:refinement} with fixed
$\alpha=100$ and the stronger universal control $\alpha=64$.

\subsection{Unlabelled selection and a labelled diagnostic oracle}

The six datasets are tested at 128 and 512 probes with the fixed,
sampled, FULL, enumeration and oracle controls in \cref{sec:6-4}.
The disjoint pools and 64 draws per sample size follow \cref{sec:3-2}.


All rule variants use \cref{eq:6-agreement} with $k=10$ and $\epsilon=1$.
FULL uses all calibration queries; the diagnostic oracle uses held-out
labels (\cref{tab:6-controls}).


\subsection{Full-pool calibration as the quality-first working point}

Full-pool calibration uses all 500 calibration queries once per cell.
Its selected arms meet the held-out $10^{-3}$ loss target in 12/12 cells;
the largest observed loss is 0.0006. Their rates are 0.908--1.000 of the
GT oracle's, and calibration costs 9.5--60.3\,ms. The geometric-mean
rate ratio to fixed 100 is $1.111\times$, or $1.193\times$ over the ten
cells where fixed 100 also meets the target. \Cref{tab:7-budget-full}
shows the complete comparison. Against fixed 64, full-pool selection is
faster in six cells, equal in two and slower in four: its main advantage
is the observed coverage of all twelve quality targets.

These cells favour FULL for quality and $S=128$ for calibration cost (\cref{sec:6-2}).


\begin{table}[tbp]
\centering
\caption[Full-pool budget calibration]{Full-pool calibration ($S=500$). Loss is relative to the best
measured held-out recall; rates are kQPS from the selected fixed arm.
``Cal.'' is one full-pool calibration call per cell, not a repeated
stream timing. Held-out counts are 500, or 499 for Stella.}
\label{tab:7-budget-full}
\footnotesize\setlength{\tabcolsep}{3pt}
% Generated by scripts/build_ch7.py; do not edit numbers manually.
\begin{tabular}{@{}lrrrrrrr@{}}
\toprule
Dataset & Probe & $\alpha_{\rm full}$ & Loss & $Q_{\rm full}$ & vs 100 & vs oracle & Cal. ms \\
\midrule
Vogue & 128 & 100 & 0.0000 & 202.9 & 1.000 & 0.908 & 10.1 \\
arXiv & 128 & 200 & 0.0000 & 151.5 & 0.728 & 1.000 & 10.0 \\
OpenAI-1536 & 128 & 8 & 0.0000 & 267.7 & 1.599 & 1.000 & 9.5 \\
OpenAI-3072 & 128 & 4 & 0.0000 & 145.0 & 1.691 & 1.000 & 17.1 \\
BGE & 128 & 64 & 0.0006 & 124.9 & 1.075 & 0.932 & 17.1 \\
Stella & 128 & 16 & 0.0006 & 110.4 & 1.153 & 1.000 & 20.6 \\
Vogue & 512 & 100 & 0.0002 & 100.1 & 1.000 & 0.969 & 20.9 \\
arXiv & 512 & 200 & 0.0000 & 69.0 & 0.831 & 1.000 & 25.1 \\
OpenAI-1536 & 512 & 8 & 0.0000 & 114.9 & 1.326 & 1.000 & 21.8 \\
OpenAI-3072 & 512 & 4 & 0.0002 & 52.1 & 1.219 & 0.977 & 40.5 \\
BGE & 512 & 64 & 0.0006 & 50.8 & 1.033 & 1.000 & 42.3 \\
Stella & 512 & 16 & 0.0006 & 36.0 & 1.041 & 1.000 & 60.3 \\
\bottomrule
\end{tabular}

\end{table}

\subsection{Lower-cost sampling and its quality risk}

\Cref{tab:7-budget-quality} reports the $S=128$ result. The median selected
budget matches the oracle budget in ten of twelve cells; it differs on
arXiv at 128 probes and OpenAI-3072 at 512. Budgets span 4--200 across workloads. Median budget and
mean recall need not describe one arm. Mean recall uses printed best recall minus mean loss, approximately to four decimals.

\begin{table}[tbp]
\centering
\caption[Sampled-budget held-out quality]{Held-out quality at $S=128$. Rule $\alpha$ is the median across
64 draws; oracle $\alpha$ uses held-out labels. $\overline R_{\rm rule}$
is reconstructed mean recall. ``Losses'' counts draws with recall loss
$>10^{-3}$ relative to the best tested held-out recall; $L_{95}$ is the
logged lower-order 95th-percentile loss estimate. $R_{100}$ is
fixed-budget recall.}
\label{tab:7-budget-quality}
\small\setlength{\tabcolsep}{3pt}
% Generated by scripts/build_ch7.py; do not edit numbers manually.
\begin{tabular}{@{}lrrrrrr@{}}
\toprule
Dataset & Probe & $\alpha$: rule/oracle & $R_{100}$ & $\overline{R}_{\rm rule}$ & Losses & $L_{95}$ \\
\midrule
Vogue & 128 & 64 / 64 & 0.9680 & 0.9677 & 5/64 & 0.0022 \\
arXiv & 128 & 100 / 200 & 0.9748 & 0.9751 & 40/64 & 0.0024 \\
OpenAI-1536 & 128 & 8 / 8 & 0.9280 & 0.9279 & 4/64 & 0.0000 \\
OpenAI-3072 & 128 & 4 / 4 & 0.9252 & 0.9252 & 0/64 & 0.0000 \\
BGE & 128 & 32 / 32 & 0.9198 & 0.9190 & 1/64 & 0.0008 \\
Stella & 128 & 16 / 16 & 0.9892 & 0.9884 & 6/64 & 0.0030 \\
Vogue & 512 & 64 / 64 & 0.9916 & 0.9915 & 2/64 & 0.0002 \\
arXiv & 512 & 200 / 200 & 0.9878 & 0.9885 & 32/64 & 0.0016 \\
OpenAI-1536 & 512 & 8 / 8 & 0.9740 & 0.9738 & 4/64 & 0.0000 \\
OpenAI-3072 & 512 & 4 / 8 & 0.9704 & 0.9702 & 0/64 & 0.0002 \\
BGE & 512 & 64 / 64 & 0.9670 & 0.9661 & 24/64 & 0.0014 \\
Stella & 512 & 16 / 16 & 0.9954 & 0.9945 & 7/64 & 0.0030 \\
\bottomrule
\end{tabular}

\end{table}

ArXiv has 40/64 and 32/64 loss crossings at 128/512 probes, despite mean loss 0.0009. Fixed 100 loses 0.0012/0.0016: better than default need not meet the target. BGE/512 has 24 crossings and $L_{95}=0.0014$; Stella reaches 0.0030; OpenAI-3072 has none (\cref{tab:7-budget-quality}). Four losses can coexist with reported $L_{95}=0$ (\cref{sec:a-budget-contract}).

\begin{figure}[tbp]
\centering
\includegraphics[width=0.82\linewidth]{ch7/budget_risk.pdf}
\caption[Held-out risk across sample sizes]{Held-out risk across sample sizes. The vertical axis is the
percentage of calibration draws losing more than $10^{-3}$ held-out recall
relative to the best tested budget. Increasing $S$ also tightens the allowed
slot-disagreement fraction because \cref{eq:6-agreement} keeps threshold one.}
\label{fig:7-budget-risk}
\end{figure}





\subsection{Estimated throughput and calibration economics}

With \cref{sec:a-budget-contract}'s mean selected-arm estimate, $S=128$ beats fixed 100 in ten cells ($1.028$--$1.691\times$). ArXiv gives 0.908/0.917 with better mean recall: no throughput payback.

Against fixed 64, sampling wins 9/12 cells: geometric mean $1.054\times$, range $0.821$--$1.367\times$ (OpenAI $1.052$--$1.367\times$). Both fixed controls fail arXiv quality, preventing a uniformly quality-matched comparison.

\begin{table}[tbp]
\centering
\caption[Sampled-calibration economics]{Calibration economics, $S=128$. Rates are kQPS and rule rate is
mean selected-arm QPS. Ratios compare with fixed 100 and fixed 64.
The ``Batches'' column gives payback relative to fixed 100 using 500
queries, or 499 for Stella; ``Never'' denotes no throughput payback.
The quality column separately marks whether fixed 100 meets the
held-out loss target.}
\label{tab:7-budget-cost}
\footnotesize\setlength{\tabcolsep}{3pt}
% Generated by scripts/build_ch7.py; do not edit numbers manually.
\begin{tabular}{@{}lrrrrrr@{}}
\toprule
Dataset & Probe & $Q_{100}$ & $\overline Q_{\rm rule}$ & Ratio & Cal. ms & Batches \\
\midrule
Vogue & 128 & 202.9 & 224.4 & 1.106 & 4.0 & 17.0 \\
arXiv & 128 & 207.9 & 188.8 & 0.908 & 4.2 & Never \\
OpenAI-1536 & 128 & 167.5 & 268.5 & 1.603 & 3.9 & 3.5 \\
OpenAI-3072 & 128 & 85.7 & 145.0 & 1.691 & 5.7 & 2.4 \\
BGE & 128 & 116.2 & 131.2 & 1.130 & 6.5 & 13.2 \\
Stella & 128 & 95.7 & 110.3 & 1.153 & 7.5 & 10.8 \\
Vogue & 512 & 100.1 & 102.9 & 1.028 & 7.5 & 55.4 \\
arXiv & 512 & 83.1 & 76.2 & 0.917 & 9.7 & Never \\
OpenAI-1536 & 512 & 86.6 & 115.0 & 1.327 & 8.8 & 6.2 \\
OpenAI-3072 & 512 & 42.8 & 52.1 & 1.219 & 13.1 & 6.2 \\
BGE & 512 & 49.2 & 51.1 & 1.039 & 15.8 & 41.2 \\
Stella & 512 & 34.6 & 36.0 & 1.041 & 21.3 & 37.1 \\
\bottomrule
\end{tabular}

\end{table}

Calibration costs 3.9--21.3\,ms. Maximum positive estimated paybacks against fixed 100/64 are 0.54/2.72\,s, or FULL 1.52/2.25\,s (\cref{sec:6-5,eq:6-payback,tab:7-budget-cost}). ArXiv controls fail quality; direct reuse below measures both cost and recall.

\subsection{Complete decision traces and finite-pool risk}
\label{sec:7-budget-trace}

All 2,304 sampled and twelve FULL curves are monotone with enumeration-matched choices. Recorded subset sums reproduce 18,432 grid counts (\cref{eq:6-finite-pool}). This verifies recorded aggregates; arbitrary subset/position invariance and general distance-only nesting remain conditional (\cref{sec:6-3}).

Maximum loss is 0.0050 (arXiv/512), 0.0046 (BGE/512) at $S=128$, and 0.0294 at $S=32$. FULL stays within 0.0006.

ArXiv/128 predicted/observed crossing rates are 0.965/0.984, 0.875/0.875, 0.620/0.625 for $S=32,64,128$ (\cref{fig:7-finite-pool}). The largest of 36 gaps is BGE/512: 0.144 (0.519/0.375). Mapping predictions to GT loss uses stable counts and held-out arm quality; calibration itself is label-free.

\begin{figure}[tbp]
\centering
\includegraphics[width=\linewidth]{supplement/finite-pool.pdf}
\caption[Finite-pool risk prediction and observation]{Finite-pool prediction and observed arXiv loss crossings. Bars
show one estimated standard error over 64 independently drawn subsets
conditional on the fixed pool, not uncertainty over future workloads.}
\label{fig:7-finite-pool}
\end{figure}

Ten alternating repeats support a low-budget timing plateau, rather than general runtime monotonicity.

\subsection{Direct calibration-plus-reuse measurement}
\label{sec:7-stream-cost}

Four cells compare fixed 64/100, sampling and FULL across five rotated rounds of 500 repeated held-out batches. Timing includes gathering, calibration, budget installation and first graph setup; build/warm-up/scan calibration are excluded (\cref{tab:7-stream-cost}).

\par\medskip\noindent\begin{minipage}{\linewidth}
\centering\captionsetup{hypcap=false}
\captionof{table}[Direct costs for 500 reuse batches]{Direct total milliseconds for 500 reuse batches, averaged over
five rounds, including budget calibration for sampled and FULL policies.
Each batch has 500 queries. Quality is discussed with the timing results.}
\label{tab:7-stream-cost}
\footnotesize\setlength{\tabcolsep}{4pt}
\begin{tabular}{@{}lrrrr@{}}
\toprule
Dataset/probes & Fixed 100 & Fixed 64 & Rule $S=128$ & FULL \\
\midrule
arXiv/128 & 1224 & 1102 & 1291 & 1680 \\
OpenAI-1536/128 & 1491 & 1264 & 852 & 860 \\
OpenAI-3072/128 & 2757 & 2160 & 1769 & 1783 \\
BGE/512 & 5185 & 5018 & 4921 & 5062 \\
\bottomrule
\end{tabular}

\end{minipage}\par\medskip

OpenAI-3072/128 always selects 4: sampled/FULL totals are 1,769/1,783\,ms versus fixed 100's 2,757\,ms, all at recall 0.9252. OpenAI-1536 FULL takes 860 versus 1,491\,ms at 0.9280; sampling takes 852\,ms at 0.9276.

Direct savings start after 4/6 batches (OpenAI-1536), 4/9 (OpenAI-3072) versus fixed 100/64; fractional fits are separate (\cref{fig:7-stream-cost}). BGE sampling takes 4,921 versus 5,185\,ms at recall 0.9658/0.9670; FULL takes 5,062\,ms at 0.9664. ArXiv FULL selects slower 200 to meet quality; fixed controls fail.

\begin{figure}[tbp]
\centering
\includegraphics[width=\linewidth]{supplement/stream-cost.pdf}
\caption[Direct cumulative calibrated-reuse time]{Direct mean cumulative time, including budget calibration and
first-batch setup. Each curve replays one fixed held-out query pool.
Policies can have different recall; timings alone are not a quality-matched
ranking. Dataset/probe settings match \cref{tab:7-stream-cost}.}
\label{fig:7-stream-cost}
\end{figure}


\section{Batch Regimes, G Scheduling and Workload Skew}
\label{sec:7-6}

This section answers the performance-regime part of RQ2: when additional
scan groups supply useful parallelism, and when their extra work dominates.

\subsection{Parallelism versus repeated selection work}



At batch 8, eight groups raise throughput $2.06$--$2.76\times$: arXiv scan drops 849$\to$216\,$\mu$s, with 16.7\,$\mu$s merging. At batch 1,000 and $\alpha=200$, independent grouping instead raises scan 3,530$\to$7,652\,$\mu$s plus 2,624\,$\mu$s merging, reaching only $0.46\times$ one-block QPS. Shared bounds reduce work, without removing this regime change.

\begin{figure}[tbp]
\centering
\includegraphics[width=\linewidth]{ch7/batch_execution.pdf}
\caption[Batch-dependent execution regimes]{Execution regimes. Left: fixed shared-bound group counts and
scheduled execution relative to the lane-refined one-block baseline
(13 cells at small batches, 11 at batch 1,000). Right: paired scan and
merge profiles at batches 8 and 1,000 for arXiv budgets 100 and 200.
Hatching denotes merge time. The right panels use the independent-bound
scan experiment; they explain a cost, not the scheduler's entire speedup.}
\label{fig:7-batch}
\end{figure}

\subsection{A fixed rule evaluated on further configurations}

The scheduler follows the fill/skew rule and pseudoquery calibration
in \cref{sec:5-5}. Four configurations set the constants before the
additional configurations were measured.


\begin{table}[tbp]
\centering
\caption[Integrated execution and scheduling gains]{Integrated execution gains over v61 L16. Entries are geometric
 means of per-cell median paired speedups from five rounds. ``Held out''
 excludes the four configurations used to set the rule; it means held-out
 configurations, not the budget experiment's held-out query split.}
\label{tab:7-scheduling}
\small\setlength{\tabcolsep}{4pt}
% Generated by scripts/build_ch7.py; do not edit numbers manually.
\begin{tabular}{@{}rrrrrrr@{}}
\toprule
Batch & Cells & $g_s=2$ & $g_s=4$ & $g_s=8$ & Scheduled & Held out \\
\midrule
8 & 13 & 1.318 & 2.014 & 2.728 & 2.727 & 2.930 \\
32 & 13 & 1.323 & 1.970 & 2.337 & 2.356 & 2.560 \\
256 & 13 & 0.997 & 1.090 & 1.099 & 1.057 & 1.041 \\
1000 & 11 & 0.959 & 0.947 & 0.886 & 1.015 & 1.012 \\
\bottomrule
\end{tabular}

\end{table}

The means use thirteen small-batch configurations and eleven full-batch ones (Stella has 999 queries); held-out subsets use nine/seven (\cref{tab:7-scheduling}). Eight groups give $0.886\times$ at batch 1,000; worst regret is BGE's 1.217 at batch 256.

Vogue's CV 0.338--0.622 forces four groups at batch 1,000. Ratios span 0.895--1.175; $(n_{\rm probe},\alpha)=(16,16),(16,64),(32,64)$ give 0.955, 0.895, 0.955. CV captures dispersion, not absolute work: the same rule enables the 0.97 lead but wastes work at low probes (\cref{sec:7-2}).

Across 250 variant/configuration pairs and 234,320 query comparisons, there are no non-tie primary differences or final-ID-set differences. Tied output order can differ. These are sequential synchronous throughput tests, not open-loop tail-latency tests.

\section{Cold Construction, Resident Memory and Tested Scale}
\label{sec:7-7}

RQ1 is assessed through cold construction, resident memory and completed
large-dataset searches; the E1 memory comparison also informs RQ4.

\subsection{Cold builds and the lifecycle boundary}

Three empty-cache repeats give Vogue 298--303\,ms training and $\approx0.56$\,s total, versus the slow archived 1.33\,s. Stella takes 10.54--10.65\,s versus archived 10.40\,s (\cref{tab:7-build-memory}).

\begin{table}[tbp]
\centering
\caption[Cold construction and resident memory]{Cold JHQ construction and separately measured resident memory.
Train/add are component medians of three supplementary normal builds;
Total is the median paired sum. Memory comes
from v71's two-group, 128-probe, $\alpha=100$ runs and includes context
and workspace. ``Residual'' is the modelled residual-code share of the
measured used total, not the fraction read per query.}
\label{tab:7-build-memory}
\small\setlength{\tabcolsep}{3pt}
% Generated by scripts/build_ch7.py; do not edit numbers manually.
\begin{tabular}{@{}lrrrrrr@{}}
\toprule
Dataset & Train s & Add s & Total s & Used GiB & Raw GiB & Residual \\
\midrule
Vogue & 1.06 & 0.27 & 1.33 & 1.47 & 2.67 & 45.4\% \\
arXiv & 0.59 & 0.52 & 1.10 & 2.58 & 6.45 & 62.5\% \\
OpenAI-1536 & 0.64 & 0.38 & 1.01 & 2.39 & 5.72 & 59.7\% \\
OpenAI-3072 & 1.34 & 0.65 & 1.99 & 4.14 & 11.43 & 69.0\% \\
BGE & 3.75 & 3.67 & 7.42 & 11.93 & 38.50 & 80.7\% \\
Stella & 3.74 & 6.67 & 10.40 & 20.44 & 67.81 & 82.9\% \\
\bottomrule
\end{tabular}

\end{table}

\Cref{sec:3-3,sec:a-e1-protocol} define startup/build boundaries; \cref{fig:7-build} compares repeated builds.


\begin{figure}[tbp]
\centering
\includegraphics[width=\linewidth]{ch7/construction.pdf}
\caption[Recorded construction costs]{Recorded construction costs, with CAGRA int8 included. JHQ is
the median and range of three supplementary normal cold builds with warm
input pages; other bars are medians across successful
measured index configurations, with minimum--maximum whiskers. Int8
includes dataset quantization; RaBitQ excludes the required serialization
round trip. Missing large-dataset bars denote unattempted configurations,
not measured build failures. The vertical axis is logarithmic.}
\label{fig:7-build}
\end{figure}

Stella takes 74.55/82.45\,s with isolated CAGRA int8 (quantisation included), and median 86.59\,s across three IVF-PQ configurations, versus JHQ's 10.54--10.65\,s. Quality/resident memory are separate comparisons.

Five arXiv builds span 0.0012 recall. Stella startup takes 18.2--18.6\,s with warm pages, 129.9\,s on first read. Sampled 20\,ms peaks are BGE 26.08\,GiB and Stella 22.58\,GiB. BGE's 12\,GiB assignment buffers precede its 23.05\,GiB gather payload, so they must not be summed (\cref{eq:4-assignment-scratch}).

\subsection{Resident allocation and where compression helps}

\Cref{fig:7-memory} combines the component model in \cref{sec:4-4}
with measured resident totals. Its unattributed remainder is retained,
and the preceding construction measurements give the sampled peaks.


\begin{figure}[tbp]
\centering
\includegraphics[width=\linewidth]{ch7/memory.pdf}
\caption[Resident-memory decomposition]{Resident-memory decomposition at two-group tables, 128 probes
and $\alpha=100$. Component sizes are modelled, markers are measured used
totals, and hatching is the unattributed remainder. The totals include
context and workspace. The preceding construction measurements give sampled peaks; these
resident markers are not peaks.}
\label{fig:7-memory}
\end{figure}

Residuals occupy 80.7\% of BGE's 11.93\,GiB and 82.9\% of Stella's 20.44\,GiB totals, including context/workspace. Raw vectors need 38.50/67.81\,GiB. Smaller $\alpha$ saves reads, not stored residual bytes (\cref{sec:4-4}).



The E1 measurements provide a separate cross-system view in
\cref{tab:7-e1-memory}. Subtracting the baseline-measured 507.2\,MiB
context from JHQ's process total gives approximately 0.29--0.39 of
CAGRA fp32's index allocation. JHQ's figure includes workspace, whereas
CAGRA's is the build-time index delta. These are parameter-sweep ranges,
not memory-matched recall points or a comparison of equal peak budgets.
The CAGRA int8 allocations in the same table are of a similar order to
JHQ's: compression relative to fp32 is shared by competing designs.
The component model in \cref{tab:4-bytes} explains this comparison.

\begin{table}[tbp]
\centering
\caption[E1 allocation ranges]{E1 allocation ranges. JHQ is process total minus a 507.2\,MiB
context estimate, including query workspace; both CAGRA columns are index
allocations measured during build. Ranges cover the parameter sweep.}
\label{tab:7-e1-memory}
\footnotesize\setlength{\tabcolsep}{4pt}
% Generated by scripts/build_ch7.py; do not edit numbers manually.
\begin{tabular}{@{}lrrr@{}}
\toprule
Dataset & JHQ net process GiB & CAGRA fp32 index GiB & Ratio range \\
\midrule
Vogue & 0.90--1.01 & 2.89--2.91 & 0.31--0.35 \\
arXiv & 2.01--2.19 & 6.84--6.98 & 0.29--0.32 \\
OpenAI-1536 & 1.79--2.33 & 5.95--5.97 & 0.30--0.39 \\
OpenAI-3072 & 3.54--3.82 & 11.67--11.68 & 0.30--0.33 \\
\bottomrule
\end{tabular}

\end{table}

\subsection{Large-dataset feasibility and competitive regions}

BGE contains 10,091,524 vectors and Stella 17,776,615, both with
1,024 dimensions and 32,768 IVF lists. Stella's queries are sampled from
the corpus; its result therefore characterizes that benchmark rather
than an independently collected external-query distribution.
\Cref{fig:7-large} includes
successful CAGRA int8 and IVF-PQ measurements. Full-precision CAGRA, RaBitQ configurations and FP32 brute force were excluded when computed storage exceeded capacity. These are planned exclusions; no out-of-memory or recall failure was measured.

\begin{figure}[tbp]
\centering
\includegraphics[width=\linewidth]{ch7/large_frontiers.pdf}
\caption[Large-dataset search frontiers]{Large-dataset host-to-host frontiers, including the successful
isolated CAGRA int8 reruns. BGE/Stella use 1,000/999 queries. Endpoints mark tested-grid coverage.}
\label{fig:7-large}
\end{figure}

CAGRA int8 fits Stella. On BGE it leads $1.35$--$2.96\times$ at recall 0.88--0.95; JHQ leads $1.086\times$ at int8's maximum 0.9579. Maximum BGE/Stella recalls are int8 0.9579/0.9783, JHQ 0.9663/0.9952, PQ 0.9525/0.9896.

Stella measured points give JHQ 185.4 kQPS/0.978679 recall and int8 66.3 kQPS/0.978278. Their $2.80\times$ ratio is not recall-matched. Int8 lacks much of the upper JHQ coverage (\cref{app:reproducibility}).



\section{Numerical Robustness, Anomalies and Recall Attribution}
\label{sec:7-8}



\subsection{Quality resolution and repeated execution}

A single GT hit changes Recall@10 by
\[
 \Delta r=\frac1{nk}=
 \begin{cases}10^{-4},&n=1000,\\2\times10^{-4},&n=500.\end{cases}
\]
Stella's $n=499$ step is slightly larger. This is resolution, not uncertainty.

E1 retains 369 baseline repeat-ID flags, including tied reordering. These repeated-call checks differ from independent builds; the 401 convention/cross-protocol checks are in \cref{sec:a-e1-protocol}.

Near 0.97, unsnapped frontier slopes give 1.01--3.17\% QPS change per $10^{-3}$ recall. Five-build ranges are Vogue 0.0008, arXiv 0.0012, OpenAI-1536 0.0004, OpenAI-3072 0.0002: few-percent local sensitivity, not a retrained-frontier bound. Cached-model FP32 checks retain OpenAI's advantage (\cref{sec:7-precision}); arXiv/LUT near ties remain inconclusive.

RaBitQ's median five-round ranges are 3.79\%, 3.34\%, 2.87\%, 5.25\% of median QPS on Vogue/arXiv/OpenAI-1536/3072. These spreads are not confidence intervals. To reduce occasional 30--60\,ms calls' influence, we recompute the 1,000-query rate as $10^6/\operatorname{median}(\text{call ms})$.
At recall 0.97, the resulting JHQ/RaBitQ ratios are 1.063, 0.995,
1.118 and 1.100. The three favourable directions and arXiv near parity
remain, so the result is not solely a consequence of long calls.

\subsection{Exactness is attached to a particular stage}
\label{sec:7-recall-attribution}

The stage contracts in \cref{tab:5-query-stages} distinguish computed-score
selection from raw-space kNN; lane rounding, cutoff ties and process
variation qualify different stages.


The archived reconstruction pilot measures the remaining gap to
full-probe reconstructed scoring.
Across seven workload--budget arms of 128 queries each, scoring all
probed candidates in FP64 improves recall over GPU outputs at selected budgets by
0--0.002344. For arXiv at 128 probes, $\alpha=200$ matches the full-IVF
reconstructed answer set on all 128 queries. At 100, the recall is
0.9734375 versus 0.9750 for full-IVF scoring: an additional 0.0015625
budget deficit beside the joint routing/reconstruction deficit of
0.0250 (\cref{sec:6-2}). These pilot aggregates include two budgets on one dataset--probe workload. FP64/GPU arithmetic differences prevent pure budget attribution in every arm. Separating routing/representation needs stage-wise GT coverage.
